% !PRINTED 520
\weberpage{520}
\gesperrt{1. Bei einer irreducibeln imprimitiven Gleichung
k\"onnen in den einzelnen Reihen die Wurzeln
so bezeichnet werden, dass die Theilgleichungen
$f(x, y_\alpha) = 0, f(x, y_\beta) = 0, \ldots f(x, y_\sigma) = 0$ im K\"orper
$\Omega(\psi)$ alle dieselbe Gruppe bekommen.}

Will man die Bezeichnung der Wurzeln nicht in der angegebenen
 Weise w\"ahlen, so werden die Gruppen wenigstens
isomorph.

Wir stellen noch die Frage, die gewissermaassen durch die
Umkehrung dieses Satzes beantwortet wird: \gesperrt{Wann hat eine
transitive Gruppe einen intransitiven Normaltheiler?}

Ist $P$ eine transitive Permutationsgruppe und $Q$ ein intransitiver
 Normaltheiler von $P$, so m\"oge
\[
A = \alpha, \alpha_1 \ldots \alpha_{r-1}
\]
eines der Systeme sein, deren Elemente durch $Q$ nur unter einander
 vertauscht werden, so jedoch, dass die Elemente von $A$
durch $Q$ transitiv verbunden sind, d.\,h. dass durch Permutationen
aus $Q$ jedes $\alpha$ in jedes andere $\alpha$ \"ubergehen kann. Wenn nun
durch ein Element $\pi$ aus $P$ die Gr\"ossen $A$ in
\[
B = \beta, \beta_1 \ldots \beta_{r-1}
\]
\"ubergehen, so m\"ussen diese Elemente entweder alle mit den $\alpha$
\"ubereinstimmen, oder alle von den $\alpha$ verschieden sein; denn
durch die Gruppe $\pi^{-1} Q \pi$, die nach Voraussetzung mit $Q$ identisch
 ist, werden die $\beta$ nur unter einander vertauscht, und zwar
ebenso wie die $\alpha$, also transitiv. Wenn also in $B$ nur ein Theil
der $\alpha$ vork\"ame, so w\"urden gegen die Voraussetzung durch $Q$
diese $\alpha$ mit anderen nicht in $A$ vorkommenden Gr\"ossen vertauscht
 werden.

Sind durch die $A$ nicht alle Elemente ersch\"opft, auf die
sich die Permutationen von $P$ erstrecken, so giebt es, da $P$
transitiv ist, eine Permutation $\pi$, durch die $A$ in ein ganz davon
verschiedenes System $B$ \"ubergef\"uhrt wird. Ist $\varkappa$ eine Permutation
 aus $Q$, so wird durch die gleichfalls zu $Q$ geh\"orige Permutation
 $\pi^{-1} \varkappa \pi$ dieselbe Vertauschung unter den Indices der $\beta$
hervorgerufen wie durch $\varkappa$ unter den Indices der $\alpha$. Daraus
folgt, dass durch $Q$ auch die $\beta$ nur unter einander permutirt
werden. Wenn mit $A$ und $B$ die Elemente noch nicht ersch\"opft
sind, so kann man in derselben Weise ein drittes System $C$
bilden, das durch $\pi'$ aus $A$ und durch $\pi^{-1}\pi'$ aus $B$ entsteht,
%CONT>>
